% Generated draft: requires mathematical and source review.
\documentclass[11pt]{article}
\usepackage[margin=25mm]{geometry}
\usepackage{fontspec}
\setmainfont{Latin Modern Roman}
\usepackage{amsmath,amssymb,mathtools,unicode-math}
\setmathfont{Latin Modern Math}
\usepackage{enumitem,needspace,xurl,hyperref}
\setlength{\parindent}{0pt}
\setlength{\parskip}{5pt}
\setlength{\emergencystretch}{3em}
\newcommand{\N}{\mathbb{N}}
\newcommand{\Z}{\mathbb{Z}}
\newcommand{\Q}{\mathbb{Q}}
\newcommand{\R}{\mathbb{R}}
\newcommand{\C}{\mathbb{C}}
\newcommand{\F}{\mathbb{F}}
\newcommand{\A}{\mathbb{A}}
\newcommand{\PP}{\mathbb{P}}
\newcommand{\E}{\mathbb{E}}
\newcommand{\eps}{\varepsilon}
\newcommand{\dd}{\,\mathrm{d}}
\begin{document}
\section*{UnsolvedMath ID: 30003548}
\textbf{Source problem code:} catalogue entry
\subsection{Definitions and mathematical statement}
Let \(n\) denote the dimension; the supplied record does not specify its admissible range. For a smooth Riemannian manifold \((M,g)\), closed means compact and without boundary. Write \(\operatorname{scal}_g(p)\) for its scalar curvature at \(p\), and define
\[
\operatorname{inj}(M,g)=\inf_{p\in M}\operatorname{inj}_g(p),
\]
where \(\operatorname{inj}_g(p)\) is the supremum of the radii \(r>0\) for which the exponential map \(\exp_p\) restricts to a diffeomorphism from the open radius-\(r\) ball in \(T_pM\) onto its image. Let \(S^n=\{x\in\R^{n+1}:\lVert x\rVert=1\}\) carry its standard smooth structure.

The supplied question asks whether there exists \(\varepsilon(n)>0\), depending only on \(n\), such that every closed smooth \(n\)-dimensional Riemannian manifold \((M,g)\) satisfies
\[
\left[\operatorname{scal}_g(p)\ge n(n-1)\ \text{for every }p\in M
\quad\text{and}\quad
\operatorname{inj}(M,g)\ge\pi-\varepsilon(n)\right]
\ \Longrightarrow\ M\cong_{\mathrm{diff}}S^n.
\]
The injectivity-radius inequality includes equality. No additional curvature, volume, or derivative bounds are imposed in this formulation. The source does not explicitly require connectedness or specify a lower dimension bound; these are unresolved qualifications requiring review before the intended scope of the assertion can be fixed.
\subsection{Short English statement}
Does a scalar-curvature lower bound of \(n(n-1)\), together with injectivity radius sufficiently close to \(\pi\) from below, force a closed \(n\)-manifold to be diffeomorphic to the standard sphere, with the permitted deficit depending only on \(n\)? The intended dimension range and connectedness assumption remain unspecified.
\subsection{Sources}
\begin{itemize}
\item Supplied local UnsolvedMath registry record, permanent ID \(30003548\), titled \emph{Sphere Rigidity from Scalar Curvature and Injectivity Radius}; collection: Oberwolfach Reports Open Problems; exact source problem code: \texttt{catalogue entry}. The record supplies proposed year 2017 but names no proposer.
\item Supplied report citation: \emph{Analysis, Geometry and Topology of Positive Scalar Curvature Metrics} (2018), Oberwolfach Reports, pp. 2223--2298, DOI \texttt{10.4171/owr/2017/36}. \url{https://doi.org/10.4171/owr/2017/36}. The supplied background describes this collection as drawn from workshop reports and their open-problem discussions.
\item Supplied literature reference: Tuschmann--Wiemeler, differentiable stability. \url{https://arxiv.org/abs/1408.3006}.
\item Supplied literature reference: Richard, strengthened injectivity bounds. \url{https://arxiv.org/abs/2404.09573}.
\item Literature assessment dated exactly \texttt{2026-08-21}: the supplied curation labels the problem \texttt{partially\_solved} and reports differentiable sphere stability under additional curvature, volume, or derivative bounds, together with strengthened injectivity estimates in selected topologies. It describes the bare scalar-curvature and near-\(\pi\) injectivity-radius question as unresolved. This is supplied dated curation, not independently verified; its time-sensitive status and cited claims require independent checking.
\end{itemize}
\end{document}
